Debido a la digitalización que se está produciendo en la realización de todo tipo de trámites a través de Internet, actualmente la sociedad tiene una dependencia creciente de las redes y sistemas informáticos. En este contexto, la ciberseguridad se ha convertido en un campo esencial para todo el mundo, tanto organizaciones como particulares.

En este ámbito, los sistemas de detección de intrusiones tienen un papel importante, ya que analizan la actividad de una red o sistema con el objetivo de identificar comportamientos fuera de lo normal o maliciosos. No obstante, la evolución de los ataques hace que resulte necesario estudiar técnicas capaces de adaptarse a estos nuevos patrones de tráfico. Es aquí donde el aprendizaje automático representa una alternativa de interés, ya que permite la creación de modelos capaces de aprender a partir de datos y realizar predicciones sobre patrones que no se conocían previamente.

A partir del conjunto de datos CIC-IDS2017, se ha procesado el tráfico contenido en este y se ha realizado una ingeniería de características para complementarlo. Tras esto se limpia el conjunto de datos y se utiliza para entrenar y evaluar el clasificador Random Forest con el objetivo de clasificar las muestras como tráfico benigno o su tipo concreto de ataque, convirtiendo este problema en uno de clasificación multiclase. Finalmente se analizan la relevancia de las características definidas y el rendimiento del modelo, proponiendo dos experimentos y comparando el rendimiento en cada uno de ellos.


\palabrasclave{Sistemas de detección de intrusiones, aprendizaje automático, Random Forest, CIC-IDS2017, ingeniería de características}
